\documentclass[10pt,a4paper]{amsart}
\usepackage{amsmath,amssymb,amsthm,graphicx}
% Retained minimal preamble for the Stacks Project chapter varieties.tex.
% Verbatim excerpt of the chapter preamble: Xy-pic setup (native preamble lines 14-18), the native theorem
% declarations with the mathematical macros (native preamble lines 154-217), and the project own no-op
% definitions of the reference/slogan/history display environments (native wrapper.tex lines 1-3) that the
% counted statement uses for its slogan. Copied mechanically; nothing added, removed or redefined.
% For commutative diagrams we use Xy-pic
\usepackage[all]{xy}

% We use 2cell for 2-commutative diagrams.
\xyoption{2cell}
% Theorem environments.
%
\theoremstyle{plain}
\newtheorem{theorem}[subsection]{Theorem}
\newtheorem{proposition}[subsection]{Proposition}
\newtheorem{lemma}[subsection]{Lemma}

\theoremstyle{definition}
\newtheorem{definition}[subsection]{Definition}
\newtheorem{example}[subsection]{Example}
\newtheorem{exercise}[subsection]{Exercise}
\newtheorem{situation}[subsection]{Situation}

\theoremstyle{remark}
\newtheorem{remark}[subsection]{Remark}
\newtheorem{remarks}[subsection]{Remarks}

\numberwithin{equation}{subsection}

% Macros
%
\def\lim{\mathop{\mathrm{lim}}\nolimits}
\def\colim{\mathop{\mathrm{colim}}\nolimits}
\def\Spec{\mathop{\mathrm{Spec}}}
\def\Hom{\mathop{\mathrm{Hom}}\nolimits}
\def\Ext{\mathop{\mathrm{Ext}}\nolimits}
\def\SheafHom{\mathop{\mathcal{H}\!\mathit{om}}\nolimits}
\def\SheafExt{\mathop{\mathcal{E}\!\mathit{xt}}\nolimits}
\def\Sch{\mathit{Sch}}
\def\Mor{\mathop{\mathrm{Mor}}\nolimits}
\def\Ob{\mathop{\mathrm{Ob}}\nolimits}
\def\Sh{\mathop{\mathit{Sh}}\nolimits}
\def\NL{\mathop{N\!L}\nolimits}
\def\CH{\mathop{\mathrm{CH}}\nolimits}
\def\proetale{{pro\text{-}\acute{e}tale}}
\def\etale{{\acute{e}tale}}
\def\QCoh{\mathit{QCoh}}
\def\Ker{\mathop{\mathrm{Ker}}}
\def\Im{\mathop{\mathrm{Im}}}
\def\Coker{\mathop{\mathrm{Coker}}}
\def\Coim{\mathop{\mathrm{Coim}}}

% Boxtimes
%
\DeclareMathSymbol{\boxtimes}{\mathbin}{AMSa}{"02}

%
% Macros for moduli stacks/spaces
%
\def\QCohstack{\mathcal{QC}\!\mathit{oh}}
\def\Cohstack{\mathcal{C}\!\mathit{oh}}
\def\Spacesstack{\mathcal{S}\!\mathit{paces}}
\def\Quotfunctor{\mathrm{Quot}}
\def\Hilbfunctor{\mathrm{Hilb}}
\def\Curvesstack{\mathcal{C}\!\mathit{urves}}
\def\Polarizedstack{\mathcal{P}\!\mathit{olarized}}
\def\Complexesstack{\mathcal{C}\!\mathit{omplexes}}
% \Pic is the operator that assigns to X its picard group, usage \Pic(X)
% \Picardstack_{X/B} denotes the Picard stack of X over B
% \Picardfunctor_{X/B} denotes the Picard functor of X over B
\def\Pic{\mathop{\mathrm{Pic}}\nolimits}
\def\Picardstack{\mathcal{P}\!\mathit{ic}}
\def\Picardfunctor{\mathrm{Pic}}
\def\Deformationcategory{\mathcal{D}\!\mathit{ef}}
\newenvironment{reference}{}{}
\newenvironment{slogan}{}{}
\newenvironment{history}{}{}
\usepackage{cleveref}
\crefname{originallemma}{original lemma}{original lemmas}
\Crefname{originallemma}{Original lemma}{Original lemmas}
\makeatletter
\newlabel{lemma-chi-tensor-finite}{{0AYT}{source}{}{sourcepointer-1}{}}
\newlabel{lemma-chi-tensor-finite@cref}{{[originallemma][0][]0AYT}{[1][0][]source}{}{}{}}
\makeatother
\crefname{originalreference}{original reference}{original references}
\Crefname{originalreference}{Original reference}{Original references}
\makeatletter
\newlabel{divisors-lemma-remove-embedded-points}{{02OL}{source}{}{sourcepointer-2}{}}
\newlabel{divisors-lemma-remove-embedded-points@cref}{{[originalreference][0][]02OL}{[1][0][]source}{}{}{}}
\makeatother
\crefname{originallemma}{original lemma}{original lemmas}
\Crefname{originallemma}{Original lemma}{Original lemmas}
\makeatletter
\newlabel{lemma-euler-characteristic-additive}{{08AA}{source}{}{sourcepointer-3}{}}
\newlabel{lemma-euler-characteristic-additive@cref}{{[originallemma][0][]08AA}{[1][0][]source}{}{}{}}
\makeatother
\crefname{originalreference}{original reference}{original references}
\Crefname{originalreference}{Original reference}{Original references}
\makeatletter
\newlabel{coherent-lemma-relative-affine-cohomology}{{089W}{source}{}{sourcepointer-4}{}}
\newlabel{coherent-lemma-relative-affine-cohomology@cref}{{[originalreference][0][]089W}{[1][0][]source}{}{}{}}
\makeatother
\crefname{originalreference}{original reference}{original references}
\Crefname{originalreference}{Original reference}{Original references}
\makeatletter
\newlabel{divisors-lemma-no-embedded-points-endos}{{02OM}{source}{}{sourcepointer-5}{}}
\newlabel{divisors-lemma-no-embedded-points-endos@cref}{{[originalreference][0][]02OM}{[1][0][]source}{}{}{}}
\makeatother
\crefname{originalreference}{original reference}{original references}
\Crefname{originalreference}{Original reference}{Original references}
\makeatletter
\newlabel{divisors-lemma-regular-meromorphic-section-exists}{{02OZ}{source}{}{sourcepointer-6}{}}
\newlabel{divisors-lemma-regular-meromorphic-section-exists@cref}{{[originalreference][0][]02OZ}{[1][0][]source}{}{}{}}
\makeatother
\crefname{originalreference}{original reference}{original references}
\Crefname{originalreference}{Original reference}{Original references}
\makeatletter
\newlabel{divisors-lemma-make-maps-regular-section}{{02P2}{source}{}{sourcepointer-7}{}}
\newlabel{divisors-lemma-make-maps-regular-section@cref}{{[originalreference][0][]02P2}{[1][0][]source}{}{}{}}
\makeatother
\title{Euler characteristic of twists is a numerical polynomial of degree at most the support dimension}
\author{The Stacks Project Authors}
\date{stacks-project:varieties.tex:local-snapshot-20261011; GNU FDL 1.2 or later}
\begin{document}
\maketitle
\noindent\textit{Editorial note.} All delivered mathematical text is the author's own, copied verbatim from the pinned Stacks Project chapter file varieties.tex; no statement or proof step was rewritten, reordered or removed.\par\medskip
\noindent\textit{Editorial note.} Dependency disclosure: the closing step 'A simple arithmetic argument shows' summarizes the standard finite-difference characterization of numerical polynomials; every preceding reduction is fully spelled in the span.\par\medskip
\noindent\textit{Allowed standard prerequisites (editorial).}\begin{itemize}
\item delivered companions: additivity and chi on finite support. delivered companions lemma-euler-characteristic-additive, lemma-chi-tensor-finite. Additivity in short exact sequences; constant base case.
\item Divisors toolbox (statement level, cross-chapter). Divisors, Lemmas divisors-lemma-remove-embedded-points, divisors-lemma-no-embedded-points-endos, divisors-lemma-regular-meromorphic-section-exists, divisors-lemma-make-maps-regular-section. Embedded-point removal, support normalization, regular meromorphic sections and the regular maps.
\item Cohomology of Schemes, affine-hypercover cohomology agreement. Cohomology of Schemes, Lemma coherent-lemma-relative-affine-cohomology. chi agreement along closed immersions.
\end{itemize}\medskip
\ifdefined\hypertarget\hypertarget{sourcepointer-1}{}\fi
\noindent\textit{Source reference.} Original lemma 0AYT: sources/stacks-chow/tags.txt, source lines 13801--13801. Cross-chapter standard prerequisite cited by the author inside the counted span; it is used as a pointer and its own proof is not part of this item.\par\medskip
\ifdefined\hypertarget\hypertarget{sourcepointer-2}{}\fi
\noindent\textit{Source reference.} Original reference 02OL: sources/stacks-chow/tags.txt, source lines 3520--3520. Cross-chapter standard prerequisite cited by the author inside the counted span; it is used as a pointer and its own proof is not part of this item.\par\medskip
\ifdefined\hypertarget\hypertarget{sourcepointer-3}{}\fi
\noindent\textit{Source reference.} Original lemma 08AA: sources/stacks-chow/tags.txt, source lines 10525--10525. Cross-chapter standard prerequisite cited by the author inside the counted span; it is used as a pointer and its own proof is not part of this item.\par\medskip
\ifdefined\hypertarget\hypertarget{sourcepointer-4}{}\fi
\noindent\textit{Source reference.} Original reference 089W: sources/stacks-chow/tags.txt, source lines 10511--10511. Cross-chapter standard prerequisite cited by the author inside the counted span; it is used as a pointer and its own proof is not part of this item.\par\medskip
\ifdefined\hypertarget\hypertarget{sourcepointer-5}{}\fi
\noindent\textit{Source reference.} Original reference 02OM: sources/stacks-chow/tags.txt, source lines 3521--3521. Cross-chapter standard prerequisite cited by the author inside the counted span; it is used as a pointer and its own proof is not part of this item.\par\medskip
\ifdefined\hypertarget\hypertarget{sourcepointer-6}{}\fi
\noindent\textit{Source reference.} Original reference 02OZ: sources/stacks-chow/tags.txt, source lines 3534--3534. Cross-chapter standard prerequisite cited by the author inside the counted span; it is used as a pointer and its own proof is not part of this item.\par\medskip
\ifdefined\hypertarget\hypertarget{sourcepointer-7}{}\fi
\noindent\textit{Source reference.} Original reference 02P2: sources/stacks-chow/tags.txt, source lines 3537--3537. Cross-chapter standard prerequisite cited by the author inside the counted span; it is used as a pointer and its own proof is not part of this item.\par\medskip
\begin{lemma}
\label{lemma-numerical-polynomial-from-euler}
Let $k$ be a field. Let $X$ be a proper scheme over $k$. Let $\mathcal{F}$
be a coherent $\mathcal{O}_X$-module. Let
$\mathcal{L}_1, \ldots, \mathcal{L}_r$ be invertible $\mathcal{O}_X$-modules.
The map
$$
(n_1, \ldots, n_r) \longmapsto
\chi(X, \mathcal{F} \otimes
\mathcal{L}_1^{\otimes n_1} \otimes \ldots \otimes
\mathcal{L}_r^{\otimes n_r})
$$
is a numerical polynomial in $n_1, \ldots, n_r$ of total degree at
most the dimension of the support of $\mathcal{F}$.
\end{lemma}

\begin{proof}
We prove this by induction on $\dim(\text{Supp}(\mathcal{F}))$.
If this number is zero, then the function is constant with value
$\dim_k \Gamma(X, \mathcal{F})$ by Lemma \ref{lemma-chi-tensor-finite}.
Assume $\dim(\text{Supp}(\mathcal{F})) > 0$.

\medskip\noindent
If $\mathcal{F}$ has embedded associated points, then we can consider
the short exact sequence
$0 \to \mathcal{K} \to \mathcal{F} \to \mathcal{F}' \to 0$
constructed in Divisors, Lemma \ref{divisors-lemma-remove-embedded-points}.
Since the dimension of the support of $\mathcal{K}$ is strictly less,
the result holds for $\mathcal{K}$ by induction hypothesis and with
strictly smaller total degree.
By additivity of the Euler characteristic
(Lemma \ref{lemma-euler-characteristic-additive})
it suffices to prove the result for $\mathcal{F}'$. Thus we may assume
$\mathcal{F}$ does not have embedded associated points.

\medskip\noindent
If $i : Z \to X$ is a closed immersion and $\mathcal{F} = i_*\mathcal{G}$,
then we see that the result for $X$, $\mathcal{F}$,
$\mathcal{L}_1, \ldots, \mathcal{L}_r$ is equivalent to the result
for $Z$, $\mathcal{G}$, $i^*\mathcal{L}_1, \ldots, i^*\mathcal{L}_r$
(since the cohomologies agree, see
Cohomology of Schemes, Lemma \ref{coherent-lemma-relative-affine-cohomology}).
Applying Divisors, Lemma \ref{divisors-lemma-no-embedded-points-endos}
we may assume that $X$ has no embedded components and
$X = \text{Supp}(\mathcal{F})$.

\medskip\noindent
Pick a regular meromorphic section $s$ of $\mathcal{L}_1$, see
Divisors, Lemma \ref{divisors-lemma-regular-meromorphic-section-exists}.
Let $\mathcal{I} \subset \mathcal{O}_X$ be the ideal of
denominators of $s$ and consider the maps
$$
\mathcal{I}\mathcal{F} \to \mathcal{F},\quad
\mathcal{I}\mathcal{F} \to \mathcal{F} \otimes \mathcal{L}_1
$$
of Divisors, Lemma \ref{divisors-lemma-make-maps-regular-section}.
These are injective and have cokernels $\mathcal{Q}$, $\mathcal{Q}'$
supported on nowhere dense closed subschemes of $X = \text{Supp}(\mathcal{F})$.
Tensoring with the invertible module
$\mathcal{L}_1^{\otimes n_1} \otimes \ldots \otimes \mathcal{L}_r^{\otimes n_r}$
is exact, hence using additivity again
we see that
\begin{align*}
&\chi(X, \mathcal{F} \otimes \mathcal{L}_1^{\otimes n_1} \otimes \ldots \otimes
\mathcal{L}_r^{\otimes n_r}) -
\chi(X, \mathcal{F} \otimes \mathcal{L}_1^{\otimes n_1 + 1}
\otimes \ldots \otimes \mathcal{L}_r^{\otimes n_r}) \\
& =
\chi(\mathcal{Q} \otimes \mathcal{L}_1^{\otimes n_1} \otimes \ldots \otimes
\mathcal{L}_r^{\otimes n_r}) -
\chi(\mathcal{Q}' \otimes \mathcal{L}_1^{\otimes n_1} \otimes \ldots \otimes
\mathcal{L}_r^{\otimes n_r})
\end{align*}
Thus we see that the function $P(n_1, \ldots, n_r)$ of the lemma has
the property that
$$
P(n_1 + 1, n_2, \ldots, n_r) - P(n_1, \ldots, n_r)
$$
is a numerical polynomial of total degree $<$ the dimension
of the support of $\mathcal{F}$. Of course by symmetry the same
thing is true for
$$
P(n_1, \ldots, n_{i - 1}, n_i + 1, n_{i + 1}, \ldots, n_r)
- P(n_1, \ldots, n_r)
$$
for any $i \in \{1, \ldots, r\}$. A simple arithmetic argument shows
that $P$ is a numerical polynomial of total degree at most
$\dim(\text{Supp}(\mathcal{F}))$.
\end{proof}

\end{document}
